\title{Parameter-Efficient Orthogonal Finetuning via Factorization}

\begin{abstract}
The widespread adoption of large foundation models is evident, yet the financial and computational costs associated with training such models from the ground up are often prohibitive. Consequently, finding ways to effectively and efficiently tailor these sophisticated models for specific downstream applications has become a central concern. This work investigates Orthogonal Finetuning (OFT), a rigorous finetuning framework aimed at task adaptation. While OFT offers strong generalizability, it typically entails a substantial number of learnable parameters owing to the inherently high dimensionality of orthogonal matrices. To tackle this inefficiency, we analyze OFT through the lens of information transfer and distill several critical requirements conducive to improved parameter efficiency. Taking inspiration from the Cooley-Tukey fast Fourier transform’s ability to transmit information efficiently, we introduce an orthogonal parameterization leveraging butterfly architectures. By integrating this approach into OFT, we propose a new parameter-efficient finetuning technique termed Orthogonal Butterfly~(BOFT). BOFT generalizes OFT by including it as a special case, thereby establishing a broader framework for orthogonal finetuning. We comprehensively assess the effectiveness of BOFT across a suite of large-scale vision transformers, large language models, and text-to-image diffusion models when adapting to a diverse array of downstream visual and language tasks.
\end{abstract}

\section{Introduction}

State-of-the-art models like ChatGPT~\cite{brown2020language,chen2021evaluating} and Stable Diffusion~\cite{rombach2022high} exemplify the extraordinary generalization potential found in contemporary large foundation models. However, the soaring capabilities of these models are accompanied by an exponential growth in parameter count

(\emph{e.g.}, GPT-3 comprises nearly 175 billion parameters). This surge has rendered the prospect of training such foundation models from the beginning infeasible for most research groups. Consequently, the advancement of the field increasingly hinges on our ability to repurpose existing foundation models for new tasks, circumventing full retraining. Efficient adaptation strategies have thus become vital. There are three primary families of approaches for effective model adaptation: \emph{model finetuning}~\cite{hulora2022,zaken2022bitfit,guo2020parameter,raffel2020exploring,qiu2023controlling,zhang2022adaptive,chavan2023one}, which entails selectively updating portions of the original model’s parameters without modifying its architecture; \emph{adapter tuning}~\cite{rebuffi2017learning,houlsby2019parameter,pfeiffer2020adapterfusion,lin2020exploring,he2021towards}, where auxiliary trainable components are introduced; and \emph{prompt tuning}~\cite{li2021prefix,lester2021power}, which attaches learnable token sequences to the input. Of these strategies, model finetuning is particularly appealing due to its conceptual simplicity, practical effectiveness, and lack of additional inference overhead.

The core concept underlying model finetuning is to maintain a close resemblance between the pretrained and finetuned models according to specific metrics, thereby retaining the acquired knowledge from pretraining. Typically, existing finetuning strategies rely on using a modest learning rate—a heuristic that implicitly ensures only minimal deviations arise between the initial and updated models. Recognizing that weight matrices exhibit inherent structural properties, an alternative systematic approach aims to conserve the relational geometry among neurons, specifically preserving the pairwise angular relationships within the matrices. This perspective motivates the development of the Orthogonal Finetuning~(OFT)~\cite{qiu2023controlling} paradigm, in which all neurons in a given layer undergo transformation via a common orthogonal matrix, guaranteeing that their pairwise angles remain invariant during finetuning. While OFT has proven effective for generalization and optimization in text-to-image diffusion model finetuning~\cite{qiu2023controlling}, one drawback is the potentially prohibitive number of trainable parameters arising from the high-dimensional nature of orthogonal matrices. To mitigate this, OFT adopts a block-diagonal design, thereby significantly reducing the parameter count. Yet, this increase in efficiency comes with trade-offs—the orthogonal matrix’s sparsity structure is fixed, and each block is manipulated independently. Consequently, although the block-diagonal configuration improves parameter efficiency, it can induce unwanted inductive biases (such as diminished expressivity, limiting the ability to replicate conventional linear transformations), and, crucially, optimal strategies for determining an effective sparsity pattern remain elusive.

A central challenge in this domain involves constructing a dense orthogonal matrix that remains efficient with respect to the number of parameters. At first glance, this appears intractable since representing a dense $d$-dimensional orthogonal matrix typically demands $\mathcal{O}(d^2)$ parameters. However, our method innovatively addresses this by assembling a dense orthogonal matrix through the composition of several sparse factorized matrices. The core idea driving this strategy is the realization that parameter count can be minimized by exchanging computational overhead for reduced parameter storage. In this scheme, since the orthogonal matrix is formed as a product of sparse matrices, repeated multiplications are necessary during each training pass. To contextualize this approach, we frame the dense orthogonal matrix construction as analogous to an information transmission challenge. Within this analogy, generating a dense orthogonal matrix via products of sparse matrices is akin to relaying information across a graph with a grid structure. This information transmission paradigm facilitates the creation of diverse sparse matrix factorizations that strike a balance between a limited number of trainable parameters and sufficient expressive power to approximate dense matrices.

To improve parameter efficiency in this information transmission setting, we borrow ideas from the butterfly networks inherent to the Cooley-Tukey fast Fourier transform algorithm~\cite{cooley1965algorithm}, which efficiently exchange information between different nodes~\cite{klasing1994broadcasting}. The butterfly graph structure in the Cooley-Tukey algorithm naturally inspires a class of sparse matrix factorizations known as butterfly factorization. By employing butterfly factorization, a dense $d$-dimensional matrix can be represented as a product of $\mathcal{O}(\log d)$ sparse matrices, each containing $\mathcal{O}(d)$ nonzero elements. By constraining each sparse factor to be orthogonal, our parameterization achieves orthogonality with only $\mathcal{O}(d\log d)$ parameters—a significant reduction from the conventional approach. Building on this parameter-efficient orthogonalization, we introduce a novel finetuning paradigm termed Orthogonal Butterfly~(BOFT). BOFT generalizes OFT, embedding it as a particular instance, and offers a unifying orthogonal finetuning framework. Both block-diagonal and butterfly structures share the defining property of sparsity, which reduces the effective parameter count in orthogonal matrices by limiting the number of nonzero elements. Comparatively, our method stands alongside low-rank approaches like LoRA~\cite{hulora2022}; both low-rank and sparse structures are foundational in the broader class of structured matrices~\cite{candes2011robust} for enabling parameter-efficient models.

In contrast to the block-diagonal design utilized by OFT, which mediates between expressivity and regularity, BOFT leverages the butterfly structure to create a more continuous spectrum between matrices belonging to the complete orthogonal group (enhancing expressivity) and the identity matrix (ensuring regularity). This property enables identification of a more compact set of hypotheses within the orthogonal group that is well-suited for various downstream tasks. Since the butterfly structure underlies many efficient linear transforms—such as the discrete Fourier transform and discrete cosine transform—we posit that adopting this structured approach introduces an implicit inductive bias. Such a bias promotes better generalization and mitigates overfitting. Our reasoning is informed by the observation that the butterfly structure is capable of representing numerous foundational linear transforms, an ability that the block-diagonal arrangement in OFT fundamentally lacks. Our primary contributions are as follows:

\begin{itemize}[leftmargin=*,nosep]
    \item We approach the challenge of parameter-efficient orthogonal finetuning through a new lens, conceptualizing it as an information transmission problem across a grid-structured graph. This reframing converts the construction of parameter-efficient dense orthogonal matrices into a problem of optimizing information flow.
    \item Drawing inspiration from the butterfly structures employed in the Cooley-Tukey algorithm, we introduce Orthogonal Butterfly—an efficient orthogonal finetuning strategy—within our proposed information transmission paradigm.
    \item We offer theoretical explanations for BOFT's ability to drastically decrease the number of required trainable parameters while retaining strong expressivity and stability during optimization. Furthermore, the matrix factorization in BOFT supports a compelling form of weight interpolation (refer to Figure~\ref{fig:boft_interpolation}).
    \item For the first time, we extend orthogonal finetuning to a broad array of applications beyond controllable text-to-image generation~\cite{qiu2023controlling}, demonstrating its wide applicability as a general-purpose finetuning technique. We evaluate BOFT on tasks spanning both computer vision and natural language processing domains, consistently observing that it outperforms contemporary leading approaches, thus confirming its excellent parameter efficiency and capacity for generalization.
\end{itemize}

\textbf{Parameter-efficient finetuning (PEFT)}. With the advent of increasingly expansive and capable foundation models (\emph{e.g.}, \cite{brown2020language,radford2021learning,rombach2022high,kirillov2023segment}), the challenge of efficiently adapting these models for specific downstream tasks—while minimizing the number of additional tunable parameters—has drawn significant attention~\cite{treviso2023efficient,ding2023parameter,lialin2023scaling}. Numerous PEFT strategies have emerged~\cite{houlsby2019parameter,aghajanyan2020intrinsic,hulora2022,edalati2022krona,wang2022adamix,gheini2021cross,zaken2022bitfit,guo2020parameter,sung2021training,ansell2022composable,lester2021power,li2021prefix,vu2022spot,he2021towards,mao2021unipelt,karimi2021compacter,liu2022few,sung2022lst,chen2022parameter,jia2022visual,chen2022adaptformer,zhang2022neural,jie2023fact,lian2022scaling,luo2023towards}, with reparameterization-based techniques~\cite{aghajanyan2020intrinsic,hulora2022,edalati2022krona} forming a central area related to this study. The LoRA method~\cite{hulora2022} introduces updates to the pretrained weight matrices by superimposing a product of two low-rank matrices, and has demonstrated notable results for various natural language processing applications. Unlike LoRA’s fixed-rank approach, subsequent works~\cite{zhang2022adaptive,valipour2022dylora,zhang2023increlora} propose strategies that vary the rank across layers, optimizing parameter distribution according to each layer’s needs. The tractability and effectiveness of these low-rank decompositions have led to widespread adoption~\cite{dettmers2023qlora,zi2023delta,chavan2023one}. Building on insights from hyperspherical energy’s relevance to model generalization~\cite{liu2018learning,liu2021orthogonal}, \cite{qiu2023controlling} developed orthogonal finetuning (OFT), which leverages orthogonal transformations within each layer for text-to-image diffusion model adaptation. This technique, where each layer’s neurons are rotated by an orthogonal matrix, is associated with improved generalization and yields more stable training behavior compared to LoRA. However, OFT typically involves a larger number of trainable parameters relative to LoRA, highlighting the need for greater parameter efficiency within OFT itself. Additionally, the broader applicability of OFT to tasks beyond text-to-image diffusion (as in \cite{qiu2023controlling}) remains an open question. By introducing butterfly factorization, BOFT enhances OFT’s efficiency, thereby enabling comprehensive evaluation of orthogonal finetuning across a wider range of adaptation scenarios.

\textbf{Butterfly structures}. The radix-2 Cooley-Tukey algorithm performs the $N$-point discrete Fourier transform by recursively dividing it into pairs of $\frac{N}{2}$-point transforms, giving rise to the characteristic butterfly structure—essentially, the product of a sequence of sparse matrices, collectively referred to as a butterfly matrix. Such matrices have found use in parameterizing orthogonal matrices (circumventing the need for pivoting in Gaussian elimination and accelerating computations~\cite{parker1995random}), stabilizing recurrent neural network training~\cite{jing2017tunable}, and facilitating efficient kernel approximations~\cite{munkhoeva2018quadrature}. Additionally, \cite{dao2019learning,dao2019kaleidoscope} have exploited butterfly parameterizations for fast linear transforms, while \cite{chen2022pixelated,dao2022monarch} employed butterfly matrices to scale up neural network training efficiency. These structures have also been leveraged for high-speed matrix-vector multiplication~\cite{michielssen1996multilevel,o2010algorithm}, data-sparse matrix approximations~\cite{li2015butterfly}, and more effective network transmission protocols~\cite{klasing1994broadcasting,li2003linear,rajasingh2016transmission}. Distinct from these earlier approaches, the present work applies butterfly structures specifically to optimize the parameter efficiency of OFT.

\section{An Information Transmission View on Orthogonal Finetuning}

We begin by reviewing fundamental aspects of OFT. In order to adapt a pretrained weight matrix, OFT reformulates the target matrix as the product of a trainable orthogonal matrix and the fixed pretrained weights. Unlike LoRA, which modifies weights by adding a low-rank matrix, OFT employs a multiplicative approach with an orthogonal matrix. For parameter efficiency, LoRA leverages a low-rank factorization, whereas standard OFT utilizes a block-diagonal orthogonal configuration~\cite{qiu2023controlling}, where parameter efficiency improves as the diagonal block size decreases. A schematic comparison is illustrated in Figure~\ref{fig:comp_lora}. The underlying rationale for using orthogonal transformations in finetuning is to maintain the pairwise angular relationships between neurons~\cite{liu2018learning,liu2021orthogonal,qiu2023controlling}, thereby preserving much of the pretrained semantic content. Specifically, OFT introduces an orthogonal matrix $\bm{R}\in\mathbb{R}^{d\times d}$ for a given pretrained linear layer $\bm{W}^0\in\mathbb{R}^{d\times n}$, modifying the forward computation from $\bm{z}=(\bm{W}^0)^\top\bm{x}$ to $\bm{z}=(\bm{R}\bm{W}^0)^\top\bm{x}$, with $\bm{x}\in\mathbb{R}^d$ as input and $\bm{z}\in\mathbb{R}^n$ as output. Zero initialization is accomplished by setting $\bm{R}$ to the identity matrix at the outset. To uphold orthogonality during training, the Cayley parameterization is used as per \cite{qiu2023controlling,liu2021orthogonal}, specifically $\bm{R}=(\bm{I}+\bm{Q})(\bm{I}-\bm{Q})^{-1}$ where $\bm{Q}$ is skew-symmetric, i.e., $\bm{Q}=-\bm{Q}^\top$. For further efficiency, the orthogonal matrix $\bm{R}$ is structured in a block-diagonal manner: $\text{diag}(\bm{R}_1,\bm{R}_2,\cdots,\bm{R}_r)$, with each block $\bm{R}_i\in\mathbb{R}^{b\times b}$ for all $i$ and $br=d$. This block-diagonal approach economizes parameters but assumes that each neuron (represented as a column of $\bm{W}^0$) is partitioned into $r$ distinct groups, and dimensions in separate groups are independently transformed by different orthogonal blocks. While empirically effective, this partitioning by index does not have a principled justification for neurons, thereby highlighting the appeal of dense orthogonal matrices. Consequently, a central question emerges: \emph{Is it possible to develop a dense orthogonal matrix while retaining parameter efficiency?}

In response to this issue, we suggest constructing a dense orthogonal matrix by multiplying several sparse orthogonal matrices. To offer a cohesive yet accessible approach for analyzing the sparsity structure in orthogonal matrix factorization, we conceptualize the process of forming a dense orthogonal matrix in OFT as analogous to an information transmission scenario. Specifically, assembling a dense matrix $\bm{R}\in\mathbb{R}^{d\times d}$ from the product of $m$ square matrices, $\thickmuskip=2mu \medmuskip=2mu \bm{R} = \bm{B}_m\bm{B}_{m-1}\cdots\bm{B}_1$, can be equated to transmitting data across a grid comprised of $d\times (m+1)$ nodes, as depicted in Figure~\ref{fig:info_trans}. This information transmission analogy is inspired by the observation that a dense $d$-dimensional square matrix encapsulates fully connected pathways from one set of $d$ nodes to another. Within $\bm{R}$, a zero at $R_{ij}$ signifies that communication from the $j$-th source node to the $i$-th target node is blocked, whereas a nonzero entry allows for such transmission. Consequently, decomposing $\bm{R}$ into several matrices $\bm{B}_m\bm{B}_{m-1}\cdots\bm{B}_1$ can be viewed as a succession of information transfer steps, each governed by the connectivity encoded in $\bm{B}_i$ for every $i$. The pathway of information transmission starts according to $\bm{B}_1$ and concludes with $\bm{B}_n$. For instance, in Figure~\ref{fig:info_trans}, we illustrate the case $\bm{R}=\bm{B}_5\bm{B}_4\bm{B}_3\bm{B}_2\bm{B}_1$ by showing both the matrices’ sparsity patterns and their associated graph structure. This graph arises from expanding the matrix product. Each matrix $\bm{B}_i$ is represented as a connectivity map linking nodes at level $i$ to those at level $i+1$. More precisely, an entry $(j_1, j_2)$ in $\bm{B}_i$ indicates the presence (or absence, if zero) of a directed connection from the $j_2$-th node in the $i$-th layer to the $j_1$-th node in the $(i+1)$-th layer. For the composition $\bm{B}_5\bm{B}_4\bm{B}_3\bm{B}_2\bm{B}_1$ to result in a dense $\bm{R}$, each node at the initial level must be able to pass information to every node at the sixth level. However, considering just $\bm{R}=\bm{B}_4\bm{B}_3\bm{B}_2\bm{B}_1$, linking the nodes in levels one and five, it becomes clear that node $1$ cannot transmit to node $3$, implying that the position $(3,1)$ in the product’s sparsity pattern is zero. This same graph-structure-to-sparsity relationship is observed when examining subproducts like $\bm{R}=\bm{B}_3\bm{B}_2\bm{B}_1$ and $\bm{R}=\bm{B}_2\bm{B}_1$. In the broader context, achieving a dense $\bm{R}\in\mathbb{R}^{d\times d}$ requires that the sequence of matrices $\bm{B}_m, \dots, \bm{B}_1$ defines a network of directed edges on a $d\times(m+1)$ grid, where each edge only bridges nodes in consecutive levels (\emph{i.e.}, adjacent columns), ensuring that every origin node in the first layer can ultimately send information to every terminal node in the $(m+1)$-th layer.

The perspective of information transmission offers valuable insights into the sparsity structure characterizing the factorization matrices within OFT. As depicted in Figure~\ref{fig:blk_diag}, $\bm{R}$ exhibits a block-diagonal form in the standard OFT approach. While such a structure serves to lower the number of trainable parameters, it precludes the possibility of realizing a fully dense matrix $\bm{R}$. The central objective, then, is to assemble a dense orthogonal matrix by combining $m$ sparse orthogonal factors, all while maintaining a minimal set of trainable parameters. When considered through the lens of information transmission, two overarching design principles emerge: \emph{(i)} \textbf{dense connectivity}, where each node at the initial layer must have a pathway to every node in the terminal layer; and \emph{(ii)} \textbf{minimum free edges}, meaning that, under the orthogonality restriction, the total number of unconstrained edges should be minimized. It is important to recognize that enforcing orthogonality introduces intricate dependencies among the connections between successive layers. For instance, ensuring that each factor $\bm{B}_i$ maintains full rank—a prerequisite for orthogonality—necessitates at least $d$ edges, thus establishing a one-to-one mapping between nodes at level $i$ and level $(i=1)$. This means there must be at least $d$ edges linking adjacent layers (such as 4 in Figure~\ref{fig:info_trans}). These $d$ connections are dictated by orthogonality and are not counted among the free, trainable edges since their values are not learnable parameters (e.g., a $d\times d$ orthogonal matrix with only $d$ nonzero elements will have each of them set to $\pm 1$). Orthogonality reduces parameter count relative to unconstrained matrices, which consequently introduces further structural dependencies among the edges. For example, within a $2\times 2$ orthogonal matrix, enforcing a zero in the $(1,1)$ entry requires a zero in the $(2,2)$ entry as well; that is, the absence of one connection can necessitate the absence of another. Thus, for any particular configuration of edges, orthogonality can impose further additions or eliminations of edges. Employing the information transmission framework—by visualizing the sparsity pattern of the resulting orthogonal matrix—is especially advantageous in the context of OFT, where the focus lies exclusively on the presence of nonzero trainable elements in $\bm{R}$, with their specific values being immaterial.

A straightforward approach to densely connect two layers requires $\mathcal{O}(d^2)$ connections (that is, implementing a single dense orthogonal matrix), amounting to $d^2-d$ connections; for example, when $d=4$, this corresponds to 12 connections. In Figure~\ref{fig:info_trans}, a sample matrix factorization is illustrated that uses only $10$ edges, which is notably fewer than those in a dense orthogonal matrix. This perspective facilitates an exploration of OFT’s parameter efficiency through graph representations, making it simple to devise effective factorizations within this paradigm. Our approach is inspired by the butterfly topology leveraged in the Cooley-Tukey algorithm~\cite{cooley1965algorithm}, which achieves efficient dense connectivity between $d$ sources and $d$ targets with just $\mathcal{O}(d\log d)$ edges. For instance, although the configuration depicted in Figure~\ref{fig:info_trans} requires $10$ edges, the butterfly topology accomplishes the same dense connections using only $8$. In the following, we present how adopting the butterfly structure enhances parameter efficiency.

\section{Orthogonal Parameterization by Butterfly Factorization}

The butterfly network was first introduced within the context of the Cooley-Tukey algorithm for performing the fast Fourier transform. In Fourier analysis, minor adjustments in the frequency space can induce sweeping changes throughout the spatial domain; this echoes the nature of our information flow problem, wherein each node at the initial layer is able to transmit to every node at the final layer. Furthermore, butterfly architectures are widely used in computer network designs~\cite{solihin2015fundamentals,li2011linear} to support highly efficient message-passing schemes. Let us suppose $k\geq 2$ is a power of $2$. The \emph{butterfly factor} $\bm{B}^F(k)$ is formulated as

\begin{equation}\label{eq:but_factor}
\begin{aligned}
    \bm{B}^F(k)=\begin{bmatrix}
    \text{diag}(\bm{d}_1) & \text{diag}(\bm{d}_2) \\
    \text{diag}(\bm{d}_3) & \text{diag}(\bm{d}_4)
    \end{bmatrix}\in\mathbb{R}^{k\times k},
\end{aligned}
\end{equation}

where each $\bm{d}_i\in\mathbb{R}^{\frac{k}{2}}$ is a vector. For dimensions $d=2^N$, the $d$-dimensional \emph{butterfly matrix} $\bm{B}(d)\in\mathbb{R}^{d\times d}$ is then recursively specified by

\begin{equation}
\begin{aligned}
    \!\!\bm{B}(d)=\tilde{\bm{B}}(d,d)\!\cdot\!\begin{bmatrix}
    \bm{B}_1(\frac{d}{2})\!\!\!\!\! & \bm{0} \\
    \bm{0}\!\!\!\!\! &\bm{B}_2(\frac{d}{2})
    \end{bmatrix}= \tilde{\bm{B}}(d,d)\tilde{\bm{B}}(d,\frac{d}{2})\cdots\tilde{\bm{B}}(d,2),
\end{aligned}
\end{equation}

In this context, $\bm{B}_1(\frac{d}{2})$ and $\bm{B}_2(\frac{d}{2})$ denote butterfly matrices, each with dimensionality $\frac{d}{2}$. We introduce the \emph{butterfly component} as $\thickmuskip=2mu \medmuskip=2mu \tilde{\bm{B}}(d,k)=\text{diag}(\bm{B}^F_1(k),\cdots,\bm{B}^F_{d/k}(k))$, a $d\times d$ block-diagonal matrix in which each diagonal block has size $k$, and every $\bm{B}^F_i(k),\forall i$ represents a butterfly factor as specified in Equation~\ref{eq:but_factor}. With this construction, we are now in a position to use butterfly matrices for orthogonal matrix parameterization. This requires that each multiplicative term $\tilde{\bm{B}}(d,k),\forall k$ in the butterfly matrix $\bm{B}(d)$ must itself be orthogonal. To illustrate, consider the block-diagonal matrix $\tilde{\bm{B}}(d,2)$, which consists of $2\times 2$ blocks. Ensuring orthogonality for $\tilde{\bm{B}}(d,2)$ is straightforward; by parameterizing every block using the Cayley transform (or a $2$-dimensional rotation), as outlined in \cite{liu2021orthogonal,qiu2023controlling}, we guarantee each block’s orthogonality. Importantly, the sparsity structure of any butterfly component is simply a permutation of the sparsity pattern found in $\tilde{\bm{B}}(d,2)$, enabling all butterfly components to be easily expressed as orthogonal matrices as well. This yields a computationally efficient scheme for parameterizing orthogonal matrices, leveraging compositions of multiple $2\times 2$ orthogonal matrices. Following the generalization in \cite{chen2022pixelated}, we extend this approach to define a block butterfly component $\tilde{\bm{B}}^b(d,k)$ in which every entry of $\bm{d}_i,\forall i$ becomes a $b\times b$ matrix. To enforce orthogonality on the block butterfly component $\tilde{\bm{B}}^b(d,2)$, we simply require that each $2b\times 2b$ block be parameterized to be orthogonal. The structure of all additional butterfly components $\tilde{\bm{B}}^b(d,k)$ for $k>2$ results from block-wise permutations of the sparsity pattern in $\tilde{\bm{B}}^b(d,2)$, and these can therefore be parameterized analogously to obtain orthogonal matrices. Bringing these ideas together, the forward computation in BOFT is expressed as follows:

\begin{equation}
    \begin{aligned}\nonumber
    \bm{z}=\big(\bm{R}(m,b)\cdot\bm{W}^0\big)^\top\bm{x},~~~~\text{s.t.}~\bigg{\{}\bm{R}(m,b)=\prod_{i=1}^m \tilde{\bm{B}}^b_{(d,i)}~~\&~~\underbrace{\big(\tilde{\bm{B}}^b_{(d,j)}\big)^\top\tilde{\bm{B}}^b_{(d,j)}=\tilde{\bm{B}}^b_{(d,j)}\big(\tilde{\bm{B}}^b_{(d,j)}\big)^\top\!=\bm{I}_d}_{\forall j\in [1, m]}\bigg{\}},
    \end{aligned}
\end{equation}

Here, for the sake of conciseness, we abbreviate $\tilde{\bm{B}}^b(d,2^{m-i+1})$ as $\tilde{\bm{B}}^b_{(d,i)}$, and $\bm{I}_d$ denotes the $d$-dimensional identity matrix. The orthogonal transformation $\bm{R}(m,b)\in\mathbb{R}^{d\times d}$ is built by sequentially multiplying several orthogonal butterfly matrices. For ease of reference, BOFT instantiated with $\bm{R}(m, \frac{b}{2})$ is written as BOFT($m$, $b$), where $b\geq 2$. When $m=1$, the structure BOFT($1$, $b$) coincides with the block-diagonal OFT~\cite{qiu2023controlling} having block size $b$. In the case of BOFT($1$, $d$), this matches the classical OFT~\cite{qiu2023controlling}, where the orthogonal matrix is full and unconstrained. For BOFT($\log \frac{2d}{b}$, $b$), the block butterfly matrix $\bm{B}^{\frac{b}{2}}(d)$ serves as $\bm{R}$, producing a fully dense orthogonal matrix. More broadly, a BOFT($m$, $b$) design results in $\frac{1}{2}(b-1)dm$ learnable parameters, when finetuning a linear transformation of size $d\times n$. When utilizing the butterfly configuration with $m=\log d$ and $b=2$, the parameter count reduces to $\mathcal{O}(d\log d)$. By contrast, a conventional OFT with an unrestricted dense orthogonal matrix employs $\mathcal{O}(d^2)$ parameters, whereas a block-diagonal OFT with $r$ blocks needs $\mathcal{O}(bd)$. Thus, traditional OFT is compelled to set $b=d$ to represent a dense orthogonal matrix, yet BOFT is flexible in the choice of $b$ to realize the same.

\textbf{Identity initialization for BOFT}. When undertaking finetuning, one usually initializes from the precise parameters of the pretrained model, limiting deviation after finetuning. For instance, LoRA adopts zero-initialized low-rank matrices. In the case of BOFT, initialization is performed by setting all butterfly factors to the identity, which, under the Cayley transform, is achieved by initializing the corresponding skew-symmetric matrices as zero.

\textbf{Multiplicative Dropout for BOFT}. In LoRA~\cite{hulora2022}, Dropout is introduced within the low-rank update to guard against overfitting, and the standard Dropout method~\cite{srivastava2014dropout} fits naturally for this use case. However, due to the multiplicative update mechanism of BOFT, traditional Dropout does not apply directly. To remedy this, we put forward a multiplicative Dropout specifically tailored to BOFT. Given that $\bm{R}(m,b)$ comprises $m$ butterfly factors—each of which can be represented as a block-diagonal orthogonal matrix of size $2b\times 2b$—our multiplicative Dropout randomly selects $p_1$ percent of the butterfly factors and $p_2$ percent of blocks in every factor, substituting those blocks with the identity matrix.

\section{Intriguing Insights and Discussions}

\textbf{Expressivity of BOFT}. The butterfly framework, when combined with suitable permutations, is capable of exactly reconstructing several canonical fast linear transforms~\cite{dao2019learning,dao2019kaleidoscope} (\emph{e.g.}, fast Fourier transform, Hadamard transform). Nevertheless, the extent to which our orthogonal butterfly matrix can approximate a generic orthogonal matrix is still not fully understood. To explore this, we conduct numerical experiments to approximate a random dense orthogonal matrix~\cite{anderson1987generation} of dimension $512\times 512$. Figure~\ref{fig:para_eff} presents the results averaged across 10 distinct random initializations. Here, the y-axis corresponds to the approximation error, while the x-axis indicates the count of effective trainable parameters. Curves sharing the same color represent BOFT variants with identical block sizes; the leftmost point on each corresponds to the error achieved by BOFT($1$,$b$) (\emph{i.e.}, the standard block-diagonal OFT at block size $b$). Overall, BOFT achieves higher parameter efficiency compared to OFT. As an illustration, BOFT($9$,$2$) surpasses BOFT($1$,$16$) in expressivity while employing considerably fewer parameters. Increasing $m$ and decreasing $b$ in BOFT configurations usually improves parameter efficiency; for instance, BOFT($6$,$4$) requires much fewer parameters yet produces approximation error on par with BOFT($2$,$16$). This is because the butterfly construction encodes a more structured class within the orthogonal group than the block-diagonal approach, ensuring that BOFT is provably more expressive than OFT when the block size is held constant.

\begin{theorem}[Expressivity of BOFT]\label{thm:exp_boft}
    With an equivalent block size, BOFT exhibits greater expressivity than OFT. Any orthogonal matrix of dimension $d$ can be constructed by multiplying butterfly matrices as $\bm{B}_{d-1,1}(d)\bm{B}^\top_{d-1,2}(d)\cdots\bm{B}_{1,1}(d)\bm{B}^\top_{1,2}(d)$, where all $\bm{B}_{i,j}(d)$ are butterfly matrices.
\end{theorem}

Based on Theorem~\ref{thm:exp_boft}, a straightforward extension of BOFT emerges: the orthogonal transformation can be generalized as $\thickmuskip=2mu \medmuskip=2mu \bm{R}^G(m_1,b_1,m_2,b_2,l)=\bm{R}_{l,1}(m_1,b_1)\bm{R}_{l,2}^T(m_2,b_2)\cdots\bm{R}_{1,1}(m_1,b_1)\bm{R}_{1,2}^T(m_2,b_2)$, where $\bm{R}_{i,j}^T(m,b)$ is the orthogonal matrix applied within BOFT. Notably, by setting $m_1=m_2=\log d$, $b_1=b_2=2$ and $l=d-1$, the form $\bm{R}^G(m_1,b_1,m_2,b_2,l)$ is capable of representing every orthogonal matrix, which aligns with the kaleidoscope hierarchy~\cite{dao2019kaleidoscope}. Despite this enhanced representational capacity, increased expressiveness does not necessarily result in superior finetuning outcomes; for instance, fully expressive finetuning is universally powerful in theory but often performs suboptimally in practice. Thus, balancing the richness of expressivity and the benefits of regularity remains central for ensuring robust model adaptation. BOFT expands the set of parameters available during finetuning by imposing structural priors, allowing for a more favorable balance between model expressivity and regularization.

\textbf{Spectral properties}. Orthogonal finetuning consistently exhibits superior preservation of spectral characteristics compared to LoRA, primarily because it maintains the spectral norm of the pretrained weight matrix $\bm{W}^0$ without alteration. This can be clarified through singular value decomposition: $\bm{W}^0=\bm{U}\bm{\Sigma}\bm{V}^\top$, in which $\bm{U}$ and $\bm{V}$ are orthogonal matrices and $\bm{\Sigma}$ is a diagonal matrix containing the singular values. Both OFT and BOFT update the parameters by applying an orthogonal matrix $\bm{R}$ on the left, resulting in the adapted weights $\bm{R}\bm{U}\bm{\Sigma}\bm{V}^\top$. This transformation leaves the largest singular value (i.e., the spectral norm of $\bm{W}^0$) unchanged. Preserving the spectral norm in this manner has been shown to enhance both the stability of training and the model’s ability to generalize~\cite{miyato2018spectral,yoshida2017spectral}. Additional mathematical details are discussed in Appendix~\ref{app:sn_property}.

\textbf{Orthogonal finetuning as learning bilinear similarity}. The BOFT approach may be recast in terms of bilinear similarity learning, specifically as $\bm{w}^0_i\bm{R}\bm{x}$, where $\bm{w}^0_i$ represents the $i$-th column (i.e., neuron) of $\bm{W}^0$. In this framework, BOFT aims to determine a similarity matrix $\bm{R}$ under the strong structural requirement that $\bm{R}$ remain orthogonal. This setup directly relates to concepts from distance metric learning~\cite{xing2002distance} and the study of bilinear forms~\cite{roman2005advanced}.

\textbf{Inductive bias and generalization in BOFT}. Within BOFT, the matrix $\bm{R}(m,b)$ is typically constrained to a specific structured subset of the orthogonal group, thereby limiting the model’s hypothesis space and introducing an explicit inductive bias. The particular structured bias imposed by butterfly factorization is argued to support generalization, as it encapsulates patterns common to numerous well-established linear transforms~\cite{dao2019kaleidoscope}, such as the discrete Fourier transform, discrete sine and cosine transforms, and the Hadamard transform. Furthermore, the sparsity inherent in BOFT’s matrix factorization may serve as a source of additional, implicit inductive bias~\cite{gunasekar2017implicit,li2018algorithmic,liu2019neural}.

\textbf{Comparison to butterfly-based sparse training}. Numerous studies~\cite{dao2019learning,dao2019kaleidoscope,chen2022pixelated,dao2022monarch} have investigated sparse training utilizing butterfly parameterization. These approaches generally concentrate on directly parameterizing neural network weight matrices with butterfly structures and involve training models ab initio. \cite{dao2022monarch} extends this by fine-tuning pretrained weights, initially projecting them onto a variant of butterfly matrices and then optimizing these projections for downstream applications. In contrast, BOFT introduces a fundamentally different fine-tuning mechanism, employing transformation of weights using layer-shared weight matrices.

\input{rebuttal-results}

\section{Concluding Remarks and Limitations}

This paper introduces Orthogonal Butterfly, a broadly applicable parameter-efficient fine-tuning approach for foundation models, drawing on the butterfly architecture. The central idea underlying enhanced parameter-efficiency is to express a dense orthogonal matrix as the product of several sparse orthogonal matrices. To facilitate the search for valid matrix factorizations, we put forth a graphical information transmission paradigm. This perspective reveals that the butterfly structure can effectively satisfy our goals of constructing sparse orthogonal matrix factorizations. Experimental results demonstrate the practical strength of BOFT for fine-tuning across a variety of domains, including large language models, large-scale vision models, and text-to-image generative models. Empirical evidence supports the general superiority of BOFT as a universal method for fine-tuning.

However, BOFT is not without its shortcomings. Due to the construction of the orthogonal matrix in BOFT as a product of multiple orthogonal matrices, training incurs a marginally higher runtime cost than OFT. Optimizing the training efficiency of BOFT is therefore an open area for further exploration. Nonetheless, once fine-tuning is complete, the orthogonal matrices learned by BOFT can be fused into the pretrained model, ensuring no additional latency during inference. Additionally, it remains uncertain whether the butterfly topology represents the most optimal structure for information propagation. The proposed information transmission paradigm encourages investigation inspired by another field—computer networking—which systematically examines the efficiency of various network topologies for information delivery. This interdisciplinary perspective could foster the development of even more efficient architectures for composing dense orthogonal matrices.